\scanpage{120}
% archive foliation: 119
La donnée d'un isom $\kappa:\underline t_E\simeq\mathcal O_t(1)$ équivaut à celui de
\[
\check\kappa:\check{\underline t}_E\simeq
\underline{\mathcal O}_t(1)^\vee\simeq\underline{\mathcal O}_t(-1)
\xrightarrow{\text{via }\omega}
\underline{\mathcal O}_t(-1)(\underline{\omega})\xrightarrow{\sim}F_t
\]
ou encore, à celui d'un \deleted{ce qui \unsure{implique un épinglage}} monomorphisme
\diagram{\begin{tikzcd}[ampersand replacement=\&]
\check t_E\arrow[r,hook]\&V_J(k)\arrow[r,hook]\&k^J
\end{tikzcd}}
dont l'image soit $F_t$, ou enfin à celui d'un élément non nul de
\[
\deleted{\unsure{\delta}}\quad
V_J(\underline t_E)=\deleted{\unsure{\underline t_E^J}}\quad
\operatorname{Ker}(\underline t_E^J\xrightarrow[\text{somme}]{}\underline t_E)
\]
donc on peut considérer que la donnée de $\kappa$ revient à celle d'une famille
\[
\tag{82}
(\delta_i)_{i\in J}\in\underline t_E^J
\]
satisfaisant à certaines relations, savoir
\[
\tag{83}
\left\{\begin{array}{ll}
\text{a)}&\sum\delta_i=0\quad\text{exprimant }\delta_i\in V_J(\underline t_E)\\[3pt]
\text{b)}&\text{Relation de proportionnalité, exprimant que}\\
&\text{l'image de }\unsure{\check{\underline t}_E|\underline t_E}\hookrightarrow V_J(k)\text{ défini par }(\delta_i)\text{ est }F_t,\\
&\text{de sorte que }(\delta_i)\text{ définit }
\kappa:\underline t_E\xrightarrow{\sim}\check F_t
\underset{\text{via }\omega}{\simeq}\mathcal O_t(1)\\[3pt]
\text{c)}&\text{Relation de l'épinglage de Legendre, exprimant}\\
&\text{que }\kappa^{\otimes2}\text{ est un isom. }\unsure{\text{standart}|\text{standard}}\text{ déjà connu.}
\end{array}\right.
\]
Nous allons expliciter les $\delta_i$ en termes du choix d'un $\delta\in\underline t_E^*$, \inserted{en plus du choix de $\omega\in\underline\omega$,} comme
\[
\tag{84}
\delta_i=\varphi_i^\omega(\delta)\delta
\qquad
\begin{gathered}
\text{où }\varphi_i^\omega(\delta)\in k^*\\
(i\in J,\ \omega\in\underline\omega(J),\ \delta\in\underline t_E^*)
\end{gathered}
\]
La condition d'invariance de $\delta_i$ p.r. au choix de $\delta$ s'exprime par la relation
\[
\tag{85}
\varphi_i^\omega(\lambda\delta)=\frac1\lambda\varphi_i^\omega(\delta)
\]
